\documentclass[11.5pt]{sig-alternate}
\usepackage{hyperref}
\usepackage{tabularx}
\usepackage{graphicx}
\usepackage{blindtext}
\usepackage[utf8]{inputenc}
\usepackage[english]{babel}
\usepackage{lastpage}
\usepackage{comment}
\usepackage{dirtytalk}
\usepackage{xcolor}
\usepackage{hanging}
\usepackage{wrapfig}
\usepackage[backend=biber,style=apa]{biblatex}
\addbibresource{notation.bib}
\usepackage{authblk}
\usepackage{caption}
\usepackage{longtable}
\usepackage{graphicx,subfigure}
\usepackage{authblk}
\usepackage{enumitem}
\usepackage[utf8]{inputenc}
\usepackage{cuted}
\usepackage{fancyhdr}
\usepackage{makecell} 
\pagestyle{fancy}
\usepackage{lipsum}
\renewcommand{\headrulewidth}{0pt}
\renewcommand{\footrulewidth}{0pt}
\setlength\headheight{80.0pt}
\addtolength{\textheight}{-80.0pt}
\usepackage{amssymb}
\setcounter{tocdepth}{2}
\chead{%
  \ifcase\value{page}
  % empty test for page = 0
  \or \includegraphics[width=\textwidth]{HTML_JournalBannerIslandConference-02.png}% page = 1
  \or \includegraphics[width=\textwidth]{HTML_JournalBannerIslandConference-02.png}% page = 2
  \or \includegraphics[width=\textwidth]{HTML_JournalBannerIslandConference-02.png}% page = 3
  \or \includegraphics[width=\textwidth]{HTML_JournalBannerIslandConference-02.png}% page = 4
  \or \includegraphics[width=\textwidth]{HTML_JournalBannerIslandConference-02.png}% page = 5
  \else
  \includegraphics[width=\textwidth]{HTML_JournalBannerIslandConference-02.png}
  \fi
}
%\chead{\includegraphics[width=\textwidth]{HTML_JournalBannerIslandConference-02.png}}
\fancyfoot[LE,LO]{Overview of the proceedings of the 2022 Inclusion in Science, Learning a New Direction, \\Conference on Disability (ISLAND)\\           
DOI: 10.14448/jsesd.15.0003}
\fancyfoot[CE,CO]{{ }}
\fancyfoot[RE,RO]{\thepage}
\pagenumbering{arabic}
\hypersetup{
    colorlinks=true,
    urlcolor=blue
}
 
\let\oldabstract\abstract
\let\oldendabstract\endabstract
\makeatletter
\renewenvironment{abstract}
{\renewenvironment{quotation}%
               {\list{}{\addtolength{\leftmargin}{1em} % change this value to add or remove length to the the default
                        \listparindent 1.5em%
                        \itemindent    \listparindent%
                        \rightmargin   \leftmargin%
                        \parsep        \z@ \@plus\p@}%
                \item\relax}%
               {\endlist}%
\oldabstract}
{\oldendabstract}
\makeatother

% Left align captions
\captionsetup{justification   = raggedright,
              singlelinecheck = false}
       
    \begin{document}
\title{Overview of the proceedings of the 2022 Inclusion in Science, Learning a New Direction, Conference on Disability (ISLAND)}

\author[1]{\large \color{blue} Cary A. Supalo}
\author[2]{\large \color{blue} Jasodhara Bhattacharya}
\author[2]{\large \color{blue} Daniel Steinberg}


\affil[1]{Educational Testing Service}
\affil[2]{Princeton University}

\toappear{}

\maketitle

%% ABSTRACT

\begin{@twocolumnfalse} 

\begin{abstract}
\item 
     \textit{The 13th annual Inclusion in Science, Learning a New Direction, Conference on Disability was hosted by the Princeton Center for Complex Materials (PCCM), a National Science Foundation funded Materials Research Science and Engineering Center (MRSEC), and Princeton University on September 16-17, 2022 at Bowen Hall. This annual conference included presentations that featured innovative research done by science educators in formal and informal education contexts, ranging from pre-K-12 to higher education, and science education researchers, access technology developers, and others interested in the full inclusion of persons with disabilities into the Science, Technology, Engineering, and Mathematics (STEM) workforce. The 2022 ISLAND conference featured seventeen different presentations over the two-day period. The following is intended to give the reader an overview of the presentations that were delivered. This is not intended to be a complete summary of all aspects of the presentations that were discussed.}
    

     
\end{abstract}

\end{@twocolumnfalse}
%% AUTHOR INFORMATION

\textbf{*Corresponding Author, Cary A. Supalo} \\
\href{mailto:mlohmann@ccu.edu}{(csupalo@ets.org)} \\
\textit{Submitted March 14, 2023}\\
\textit{Accepted April 04, 2023} \\
\textit{Published online March 22, 2023} \\
\textit{DOI: 10.14448/jsesd.15.0003} \\
\pagebreak
\pagebreak

\vspace{5mm}
\section*{\vspace{140mm}}
\section*{Introduction}
\begin{large}

\textbf{The first set of three presentations occurred on the evening of Friday September 16th}, and highlighted the impact of both transformative access technologies and a supportive ecology of education on the survival of blind scientists in the STEM pipeline. 

\textbf{The first presentation on Friday evening was titled, “Changing the Landscape of STEM Education for B/VI Students,”} delivered by Dave Schleppenbach who is the founder of Tactile Engineering based in Lafayette, Indiana. This presentation discussed one of the first multi-line refreshable braille display tools that is capable of both multi-line refreshable braille display and dynamic tactile graphics all on the same screen. This technology can receive digital content and display it in real-time on the multi-line device. This technology has been branded, “Cadence,” and features a 4 line by 12 braille cell footprint. It is intended to be used in pairs to provide a 4 line by 24 cell display, or can be used as a quad which would require 4 Cadence devices that would provide an eight line by 24 cell display. The devices are Bluetooth connected and currently have a word processor, maze drawing game, and scientific calculator applications. Tactile Engineering is collaborating with Independence Science to develop a scientific data collection application that would leverage the Cadence’s tactile graphics display capability as part of a new scientific data logger device. Over all, the Cadence once it reaches commercialization has the potential to be a game changer in the blindness field. The second presentation was titled, \textbf{“High-Power, Low-Cost, Tactile Actuation Arrays Enabled by Soft Micro-cumbustion,” and was delivered by} Ronald He-isser from Cornell University in Ithaca, New York. This presentation discussed another refreshable braille technology that leverages controlled combustion to fire actuators that would enable braille dots to be raised or lowered. Although the researchers admit this is a non-conve-ntion approach to controlling refreshable braille applications, this unique approach may prove to have other commercial uses.

The final presentation on the opening evening of the 2022 ISLAND conference was titled, \textbf{“An Autobiography on Changes of Blindness,”} and was presented by Charis Glatthar who is currently an under graduate student at Metro State University, Denver located in Denver, Colorado. This presentation was a biographical account of how Charis Glatthar over came a number of her accessibility challenges as a blind student enrolled in an environmental science degree program. Her use of the Sci-Voice Talking LabQuest, along with a very progressive faculty at Metro State that encouraged her to engage in a hands-on way in her laboratory science courses as a blind student. It is this environmental encouragement that helped to facilitate Ms. Glatthar’s love for science. The evening session closed with networking time for those on-line and physically present on the Princeton University campus.

\textbf{The 2022 ISLAND conference resumed on the morning of Saturday, September 17, 2022.} The first block of three presentations focused on technological interventions, whether it is the impact of access to multi-sensory laboratory learning technologies for populations with other disabilities, or the transformative potential of ground-breaking real-time tactile graphics access for the blind in STEM (science, technology, engineering and mathematics), or sonification-based technologies for specific STEM laboratory learning applications. 

The first presentation was titled, \textbf{“B/VI Accessibility Technology Adapted for Chemistry Students with Language-Based Disabilities,”} presented by Dr. Christin Monroe from Landmark College based in Putney, Vermont. This presentation discussed how neuro diverse under graduate chemistry students used the Sci-Voice Talking LabQuest in several under graduate first year chemistry lab activities. This investigation tried to see how text-to-speech scaffoldings would positively or negatively impact student learning. Feedback on the improvement of the quality of the text-to-speech and how visual indicators on the screen that would indicate on a Cartezian graph during data collection where the data was being collected would enhance student learning was documented. Although the Sci-Voice Talking LabQuest was developed by blind scientists for blind students enrolled in STEM laboratory classes, this feedback is valuable for future versions of these types of supports that will expand the user base beyond BVI learners.

\textbf{The second presentation of the morning was titled, “Synergizing Braille and Science: Real-Time Tactile Graphics Access in Science Laboratory Settings for Students Who are Blind,”} delivered by Dr. Ashley Nashleanas from Independence Science based in West Lafayette, Indiana. This presentation discussed specifically the results of the usability study that was conducted in Summer 2022 of the Tactile Engineering Cadence multi-line refreshable braille device with the prototype scientific data collection mobile application. All ten participants positively reported the practicality of tactile graphics display during data collection was a valuable experience. Feedback on hardware design of the Cadence was also collected and documented. Over all, the ten stake holders indicated genuine excitement for this positive scientific data collection experience and encouraged this product be pursued to commercialization.

The next presentation prior to the morning break was titled, \textbf{“APH Presents the Submersible Audio Light Sensor (SALS): A Device That Fosters Inclusion in Science,”} was presented by Dr. Rosanne Hoffman from the American Printing House for the Blind in Louisville, Kentucky. This presentation featured the recently commercially released SALS device in the APH catalog and discussed how this handheld sensor interfaced with a mobile application. Once the sensor is Bluetooth connected to the mobile device, it could be used in an analogous to the original prototype SALS devices that were developed as part of the Independent Laboratory Access for the Blind or (ILAB) project that was sponsored by the National Science Foundation Research in Disabilities Education (RDE) program and hosted at Penn State University in University Park, Pennsylvania. This new SALS sensor is eligible for quota dollar purchase by k12 schools in the United States. The SALS measures color changes in test tubes and other clear containers. It uses ambient light along with the change in light transmittance within the solution to produce an audible tone. The tone decreases as the magnitude of the color change increases. The SALS can also be used for acid-base titrations, comparison of concentrations of solutions, and for chemical reactions that emit gases in the form of bubbles. The SALS can also be used to detect color changes on solid surfaces. APH has released a set of science labs that features the SALS for its customers.

\textbf{The next block of three presentations, which occurred after the morning break,} were centered on educational programming for populations with the low incidence disabilities of blindness/ low-vision and deafness/hard-of-hearing. These highlighted the need for quality accessible formal and informal educational experiences for students with these disabilities, to maintain their engagement in the STEM fields of education. 

\textbf{The first presentation after the morning break was titled, “Including BVI Students in NSF Funded Research Experience Programs,”} and was presented by Dr. Daniel Steinberg from Princeton University in Princeton, New Jersey. This presentation discussed how Princeton University was in the process of establishing a research experience for undergraduate students with disabilities with distinguished Princeton faculty. Once several faculty who are willing to host students with disabilities are identified, students with disabilities will be recruited for these opportunities. Based on the undergraduate students individualized accommodations requests are met, this summer REU experience will commence. It is the expectation of Princeton University that these students participate in a poster session upon the conclusion of their summer research experience to share their work. All students are encouraged to pursue publication of their findings.

The second presentation was titled, \textbf{“Deaf and Hard of Hearing Individuals Diversifying the STEM Workforce: through Increased Access to STEM Education,”} and was presented by Jason Nordhaus and Jessica Williams from the Rochester Institute of Technology in Rochester, New York. This presentation field tested a series of physics lesson videos with Deaf and Hard-of-hearing undergraduate students that featured ASL interpreters who were familiar with the physics content. One of the major challenges with ASL interpreters is their non-familiarity with STEM content which can make it difficult for these interpreters to provide effective ASL services. As the number r of ASL interpreters with STEM backgrounds at this time is some what limited, those that do possess this knowledgebase can be video recorded providing appropriate ASL support that can be shared with multiple students at a time. By making this experience available to DHH students increases the number of students who depend on this service to benefit from it. As this capability is applied to other STEM courses in the future, it is the hope this model can be replicated in other STEM discipline courses at RIT.

\textbf{The final presentation of the morning was titled, “An Introductory Course in Electrical Circuits and Coding for Deaf and Deaf Blind Middle School Students,”} presented by Christina Yang from the Playful Learning Lab based in St. Paul, Minnesota. The Playful Learning Lab partnered with the Metro Deaf School (MDS) to develop and deliver an Introductory class for Electronics and Computer Programming for Deaf and DeafBlind middle school students. This program featured four sets of 6-8 students went through an approximately 20-day course that met four days a week for 40 minutes per day. The content of the class focused on electrical circuits and computer programming. Some of the lessons discussed topics like conductivity, closed/open circuits, series/parallel circuits, and coding. The participating students were also introduced to the Scratch computer programming language and Makey Makey kits. The development and delivery of this course was a collaborative effort between Metro Deaf School staff and the undergraduate student researchers from the Playful Learning Lab. According to the presenters,  this class was a project-based mix of lectures, projects, and hands-on assignments that featured pedagogy in all of these areas. This was the last presentation prior to the lunch break.

\textbf{The morning sessions were followed by a keynote by Dr. Bryan Shaw, from Baylor University in Waco, Texas.} This keynote presentation, titled \textbf{"A vision quest in the chemistry lab"}, discussed how Dr. Shaw’s work that was featured in the journal, “Science Advances,” in May 2022 described how ancient technology called lithophanes could be 3-D printed today to make inclusive tactile images of scientific data for all students. The lithophanes can be used in quantifiable data analysis which has its advantages over swell form and tactile embossed graphics for blind students. Additionally, lithophanes can also have visual color thus making them usable by non-vision impaired students. This inclusive approach to data visualization illustrates how both sighted and blind students can benefit from this type of data visualization approach. Once concern about the use of 3D printed lithophanes is the extra time it takes to produce the 3D printed images. Dr. Shaw indicated many of these images take on average about 30 minutes more to produce. When you add this extra time on to the time it takes to analyze the data set within the analysis software, this work can take a number of hours. Therefore, the extra 30 minutes on the back end does not increase the overall time a significant amount.

Dr. Shaw also spent some time discussing his work with his open source mobile application, called “Cradle” or “The White Eye Detector”, which uses photographs of children’s eyes and detects abnormal reflections from the retina, which constitutes leukocoria, or “white eye”, wh-ich is a characteristic indicator of the aggressive pediatric eye cancer known as retinoblastoma, as well several other eye disorders. The impetus for this mobile application was Dr. Shaw’s lived experience, as the parent of a vision impaired son who survived retinoblastoma. It is this experience that motivated him to develop this capability that is helping parents to access a low-cost early diagnosis system to preserve the eye sight in children around the world. Cradle (which is an acronym for Computer Assisted Detector of Leukocoria) has been used by a large number of families around the world. Dr. Shaw’s motivation to continue work in the blind accessibility space are truly noteworthy and are going a long way in promoting the full inclusion of blind students in STEM education and employment.

\textbf{The seven presentations which followed the lunch-break touched on various entry-points in the larger STEM ecology, for the blind to access, survive and thrive in educational, social, and economic spaces.} These talks touched on accessible curriculum and assessments, from co-curricular programs that build scientific inquiry skills in blind STEM learners, to the importance of accessible and equitable gate-keeping assessments to make entry into STEM higher education and the workforce possible. They also highlighted the importance of teacher professional development and accessible teaching practices on student success, as well as STEM teaching-learning resources rendered accessible through the use of AI technologies. The audience was also introduced to the impact of non-profits in both raising awareness about the causes of blindness, to their role in the development of technologies to provide independent access to the infrastructure of everyday life.

\textbf{The first presentation of the afternoon was titled, “Using AI to Instantly Make STEM Images and Documents Accessible,”} and was presented by Vijayshree Vethantham from Continual Engine based in Dallas, Texas. Continual Engine’s product branded PREP and Invicta are innovations in accessible technology that enable more equitable access to STEM content. These products leverage artificial intelligence (AI), machine learning, automation, and other technologies in an innovative way to make STEM content accessible for students with print disabilities. The goal of these accessibility products are to process all types of complex STEM images at all levels of education. The PREP PDF and Document Remediation Platform is a document accessibility platform that accelerates and automates the process of making PDFs and other documents accessible for individuals with print disabilities who use screen readers. The presenters state, “PREP minimizes the need for manual intervention as it automates tagging for even complex structures like tables and lists.” Therefore, it is the hope that this new innovative way of automatically tagging STEM complex images can make large amounts of technical content accessible to students with print disabilities and thus contributing to a more inclusive accessible STEM learning experience

The second presentation of the afternoon session was titled, \textbf{“Mission Inspire: Having a BLAST with Students with Visual Impairments Learning STEM”.} Presented by Dr. Tiffany Wild from The Ohio State University in Columbus, Ohio and Dr. Tina Herzberg from the University of South Carolina Upstate , Valley Falls, South Carolina. This presentation featured a program that includes BVI learners in hands-on STEM learning experiences that encourages the use of the scientific method to solve science questions all while engaging in data collection activities. The BVI students also had the opportunity to practice their presentation skills as part of this program. The students had a culminating project to launch rockets. This program encourages BVI students to participate in state wide science fair competitions. It is the hope of the presenters that Project Inspire will serve as a model for other educators that will encourage BVI learners to more fully participate in science fair competitions.

The next presentation was titled, \textbf{“TVIs Mathematical Pedagogical Content Knowledge,”} and was delivered by Tasnim Al Shuli who is a graduate student at the University of Arizona, Tucson, Arizona. This presentation discussed a research study that investigated how teachers of the visually impaired (TVI) taught math content to students with visual impairments. This study documented teaching practices for how TVI’s communicated math content to their students with visual impairments. This study is ongoing and has not yet been completed.

\textbf{The next presentation, occurring after the afternoon break, was titled, “Information on Wolfram Syndrome,”} and was delivered by Ellie White who is the founder of the Ellie White Foundation based in  Aurora, Colorado. Wolfram Syndrome is an extremely rare genetic disorder that can be referred to as DIDMOAD (diabetes insipidus diabetes mellitus optic atrophy deafness). There is very little that is known about this disorder, however with research being conducted at Washington University in St. Louis, Missouri, new findings are being discovered that is unlocking some of this disease’s many secrets. The Ellie White Foundation’s mission is to draw awareness for families, research facilities, and medical centers. Fundraising is another goal of the foundation to support this research since little government funding is currently being invested.  It is the hope that more people will be made aware of the Wolfram Syndrome and more interest in medical research is needed to improve the quality of life for persons with this condition.

The important work of this not-for-profit was followed by \textbf{“Development, Implementation, \\and Preliminary Outcomes from the Connecting Students with Autism to Geographic Information Science \& Technology (CSA-GIST) Program,”} presented by Jamie Pearson from North Carolina State University, Raleigh, North Carolina. In this presentation, the investigators discussed the development and implementation of GIST, and highlight preliminary findings following year one of this project.  Further, the researchers know many students with autism thrive in stimulating STEM careers when given proper supports. To that end, the Connecting Students with Autism to Geographic Information Science \& Technology (CSA-GIST) study established a research-based workforce development model that increases student self-regulat-ion, interest, and motivation in Geographic Information Science\& Technology (GIST). It also expanded students understanding of GIST/STEM concepts and skill sets in an effort to prepare autistic high school students for post-secondary STEM education. Additional findings from this research will be published as they become available.

The penultimate presentation of the afternoon was titled, \textbf{“Project ATOM,”} by Dr. Martin Goldberg from Tinski Tech Inc based in New York, New York. This presentation discussed a prototype technology that was being designed for BVI persons to assist in navigation. Specifically, the last 50 feet tends to be the most difficult for blind travelers. This technology complements other commercially available visual interpreter services by providing a chaperone service for the blind individual to assist with the last fifty feet of travel.

The final presentation of the 2022 ISLAND conference was titled, \textbf{“Evaluation in a Science Assessment Context,”} and was presented by Dr. Mark Hakkinen from the Educational Testing Service, Princeton, New Jersey. This presentation featured a series of demonstrations of how audio sonification techniques could be used to sonify tabular data. Normally, BVI persons have to investigate tables cell by cell thus causing increased cognitive load concerns as tables become larger. This formal research study included seven BVI participants and investigated how these subjects could use these tools to investigate tabular data. This study also observed how these subjects could narrow down tabular data to specific points of interest. The researchers indicate more work needs to be done on the user interface and on the documentation as to how to use these prototype tools.

The 2022 ISLAND conference spanned presentations that touched on the breadth of STEM learning experiences, resources and opportunities for persons with disabilities. Some presentations demonstrated progress or the final product of prototypes unveiled at previous ISLAND conferences. Many featured techniques and approaches for teaching a range of scientific disciplines in new innovative ways. Others described and demonstrated specific methodologies that were implemented in their specific scenarios, which could be reproduced and modified for varying contexts.  More formal research conducted in the future, to determine the effectiveness of some of the methodologies presented as part of the ISLAND conference, may lead to innovative new best practices for the science education community. Multisensory science instruction is becoming more mainstream, and new innovative digital technologies are becoming more widely available to more effectively render STEM content in inclusive accessible outputs. These advances will further promote persons with disabilities entering the STEM work force. As the science education community of practice grows with individuals interested in inclusive STEM education methodologies, more discussions are necessary to advance new innovative methodologies and access technology development and implementation to advance more inclusive and equitable instruction in K-12 and higher education.  The Princeton University MRSEC outreach program, helped sponsor the program in 2022 and is planning on continuing its involvement in future ISLAND conferences.

“This conference was partially supported by NSF through the Princeton University Materials Research Science and Engineering Center DMR-1420541”
 

\end{large}
\end{document}
